The original path passes 19 native tests, 7 CPU/tester cases, and 4 XPU cases. SparseGate then passes separate synthetic-input CPU, learned-weight-derived CPU, and three-source coexistence programs. The latter executes the original GPU/NPU work before gate commands; it demonstrates coexistence on the retained path, not concurrent gate contention.

Every gate run checks all returned score/index records and every gathered byte against the independent fixture. The small-case sequence exercises FULL, valid REUSE, REINDEX, stale-epoch rejection, and FULL recovery. The stale command has no accepted DMA traffic. Native AXI waves, UCIe Flits, backend events, and memory commands are cross-checked, with binary and source digests retained. Each command counter is independently checked against its RTL busy interval, and successful DONE is checked against the final DMA write response. A separate adapter/RTL test rejects unsupported exclusive accesses and verifies ordinary-access recovery.

The waveform checker also rejects three controlled corruptions of a recorded trace: an incremented cycle count, an incremented DMA-read count, and premature DONE at the final B handshake. These overlays leave the original simulation files unchanged.

\begin{figure*}[t]
\centering\includegraphics[width=.94\textwidth]{completion_wave.pdf}
\caption{Actual native RTL VCD samples around the first successful FULL completion in the H4/N16/K4 CPU program. Levels are recorded after model evaluation; handshakes sample values stable before the rising edge. DONE commits three AXI cycles after the final B handshake. The RTL STATUS response is independently correlated with its later delivery through the return protocol path. This is a simulator waveform, not a silicon capture.}
\label{fig:completion}
\end{figure*}
\begin{table}[t]
\centering\caption{Measured RTL command counters in the learned-weight-derived H32/N64/K8 program; REINDEX has 16 candidates. R/W are DMA beat counts, unchanged across the two memory periods. Each result gathers 288 B.}
\label{tab:modes}
\small
\begin{tabular}{lrrrrr}\toprule
Mode & Scores & R & W & $1\times$ cyc. & $4\times$ cyc.\\\midrule
FULL & 64 & 360 & 80 & 55,155 & 67,969\\
REUSE & 0 & 72 & 80 & 1,949 & 5,783\\
REINDEX & 16 & 218 & 80 & 18,788 & 26,359\\
\bottomrule\end{tabular}
\end{table}
The AXI simulation period is 2 ns. Command cycles count the RTL busy interval through the final write response; host-run time also includes the protocol and software interactions. All three successful commands return eight records on this input.

Increasing the online memory period from 0.25 to 1 ns preserves the guest binary, RTL, scores, and DMA counts. All successful command counters increase; complete host finish time increases by 181.530 $\mu$s. Polling traffic can change, so this is a memory-feedback experiment rather than a fixed-instruction speedup comparison.

A separate full-shape H32/N640/K512 guest executes FULL and then REUSE to distinct destinations. It verifies all 512 score/index records, their padding, and 147,456 gathered bytes per command. FULL takes 670,572 busy cycles with 6,624 read beats and 5,120 write beats; REUSE takes 115,776 cycles with 4,608/5,120 read/write beats. This extends full-shape selection evidence through the actual online path; the small-case REINDEX/error sequence is a separate test.

\begin{figure*}[t]
\centering\includegraphics[width=.86\textwidth]{native_modes.pdf}
\caption{Measured H32/N64/K8 command latency with two online memory periods, and DMA payload (identical across periods). REUSE preserves an earlier selection; REINDEX evaluates a supplied subset. These different workloads are not interchangeable accuracy baselines. Payload includes result and gathered-KV writes, excluding host uploads, MMIO, readback, and link framing.}
\label{fig:native}
\end{figure*}
